% ==========================================================
\section*{Learning Objectives -- Mục tiêu bài học}
\addcontentsline{toc}{section}{Mục tiêu bài học}

Bài học này tiếp tục chủ đề hồi quy tuyến tính, đi sâu vào các khía cạnh nâng cao và thực hành.

Sau khi học xong, người học cần nắm được:

\begin{enumerate}[label=\arabic*.]
    \item Cách đánh giá chất lượng mô hình hồi quy.
    \item Các độ đo phổ biến: MSE, RMSE, MAE, $R^2$.
    \item Hiểu overfitting và underfitting trong hồi quy.
    \item Regularization: Ridge và Lasso regression.
    \item Feature engineering trong hồi quy.
    \item Polynomial regression.
    \item Cách xử lý outliers.
    \item Validation và cross-validation.
\end{enumerate}

% ==========================================================
\section{Model Evaluation -- Đánh giá mô hình}

\subsection{Metrics for Regression -- Các độ đo cho hồi quy}

Để đánh giá chất lượng mô hình hồi quy, ta cần các độ đo định lượng.

\subsubsection{Mean Squared Error (MSE)}

MSE là độ đo phổ biến nhất, tính trung bình bình phương sai số:

\[
    \text{MSE} = \frac{1}{n}\sum_{i=1}^{n}(y_i - \hat{y}_i)^2
\]

Trong đó:
\begin{itemize}
    \item $n$ là số mẫu;
    \item $y_i$ là giá trị thật;
    \item $\hat{y}_i$ là giá trị dự đoán.
\end{itemize}

\subsubsection{Root Mean Squared Error (RMSE)}

RMSE là căn bậc hai của MSE:

\[
    \text{RMSE} = \sqrt{\text{MSE}} = \sqrt{\frac{1}{n}\sum_{i=1}^{n}(y_i - \hat{y}_i)^2}
\]

Ưu điểm: có cùng đơn vị với biến mục tiêu.

\subsubsection{Mean Absolute Error (MAE)}

MAE tính trung bình giá trị tuyệt đối của sai số:

\[
    \text{MAE} = \frac{1}{n}\sum_{i=1}^{n}|y_i - \hat{y}_i|
\]

MAE ít nhạy cảm với outliers hơn MSE.

\subsubsection{R-squared ($R^2$)}

$R^2$ đo phần trăm phương sai được giải thích bởi mô hình:

\[
    R^2 = 1 - \frac{\sum_{i=1}^{n}(y_i - \hat{y}_i)^2}{\sum_{i=1}^{n}(y_i - \bar{y})^2}
\]

Trong đó $\bar{y}$ là giá trị trung bình của $y$.

Giá trị $R^2$:
\begin{itemize}
    \item $R^2 = 1$: mô hình hoàn hảo;
    \item $R^2 = 0$: mô hình không tốt hơn dự đoán bằng trung bình;
    \item $R^2 < 0$: mô hình tệ hơn dự đoán bằng trung bình.
\end{itemize}

% ==========================================================
\section{Regularization -- Điều chuẩn}

\subsection{Overfitting Problem -- Vấn đề quá khớp}

Khi mô hình có quá nhiều tham số, nó có thể khớp quá mức với dữ liệu huấn luyện, dẫn đến overfitting.

Regularization thêm penalty term vào hàm mất mát để hạn chế độ phức tạp của mô hình.

\subsection{Ridge Regression (L2 Regularization)}

Ridge regression thêm penalty $L_2$ norm của trọng số:

\[
    L(\mathbf{w}) = \sum_{i=1}^{n}(y_i - \hat{y}_i)^2 + \lambda \sum_{j=1}^{m}w_j^2
\]

Trong đó:
\begin{itemize}
    \item $\lambda$ là hệ số regularization;
    \item $w_j$ là các trọng số.
\end{itemize}

Ridge không đưa trọng số về 0, chỉ giảm nhỏ chúng.

\subsection{Lasso Regression (L1 Regularization)}

Lasso regression thêm penalty $L_1$ norm:

\[
    L(\mathbf{w}) = \sum_{i=1}^{n}(y_i - \hat{y}_i)^2 + \lambda \sum_{j=1}^{m}|w_j|
\]

Lasso có thể đưa một số trọng số về chính xác 0, thực hiện feature selection.

% ==========================================================
\section{Advanced Topics -- Chủ đề nâng cao}

\subsection{Polynomial Regression}

Khi quan hệ giữa biến độc lập và phụ thuộc không tuyến tính, ta có thể dùng polynomial regression:

\[
    y = w_0 + w_1x + w_2x^2 + w_3x^3 + \cdots + w_dx^d
\]

Đây vẫn là linear regression vì tuyến tính theo tham số $\mathbf{w}$.

\subsection{Feature Engineering}

Feature engineering là quá trình tạo biến mới từ biến gốc:
\begin{itemize}
    \item Tương tác giữa các biến: $x_1 \cdot x_2$;
    \item Biến đổi phi tuyến: $\log(x)$, $\sqrt{x}$, $x^2$;
    \item Binning: chia biến liên tục thành nhóm;
    \item One-hot encoding cho biến categorical.
\end{itemize}

\subsection{Handling Outliers}

Outliers có thể ảnh hưởng mạnh đến mô hình hồi quy. Cách xử lý:
\begin{itemize}
    \item Loại bỏ outliers dựa trên z-score hoặc IQR;
    \item Dùng robust regression (Huber loss);
    \item Transform dữ liệu (log transform).
\end{itemize}

% ==========================================================
\section{Model Validation -- Kiểm chứng mô hình}

\subsection{Train-Test Split}

Chia dữ liệu thành:
\begin{itemize}
    \item Training set (70-80\%): huấn luyện mô hình;
    \item Test set (20-30\%): đánh giá cuối cùng.
\end{itemize}

\subsection{Cross-Validation}

K-fold cross-validation:
\begin{enumerate}
    \item Chia dữ liệu thành $K$ fold;
    \item Mỗi lần, dùng $K-1$ fold để train, 1 fold để validate;
    \item Lặp $K$ lần, mỗi fold làm validation một lần;
    \item Tính trung bình các kết quả.
\end{enumerate}

Thường dùng $K=5$ hoặc $K=10$.

% ==========================================================
\section{Technical Glossary -- Bảng thuật ngữ}

\begin{center}
\begin{tabular}{p{0.3\textwidth}p{0.6\textwidth}}
\toprule
\textbf{Thuật ngữ} & \textbf{Giải thích} \\
\midrule
MSE & Mean Squared Error -- trung bình bình phương sai số \\
RMSE & Root Mean Squared Error -- căn MSE \\
MAE & Mean Absolute Error -- trung bình giá trị tuyệt đối sai số \\
$R^2$ & Coefficient of determination -- hệ số xác định \\
Ridge & L2 regularization \\
Lasso & L1 regularization \\
Cross-validation & Kiểm chứng chéo \\
Outlier & Điểm bất thường, ngoại lai \\
\bottomrule
\end{tabular}
\end{center}

% ==========================================================
\section{Summary -- Tóm tắt}

\begin{enumerate}
    \item Đánh giá mô hình hồi quy bằng MSE, RMSE, MAE, $R^2$.
    \item Regularization (Ridge, Lasso) giúp giảm overfitting.
    \item Ridge sử dụng L2 penalty, Lasso dùng L1 penalty.
    \item Lasso có thể thực hiện feature selection.
    \item Polynomial regression mở rộng cho quan hệ phi tuyến.
    \item Feature engineering tạo biến mới để cải thiện mô hình.
    \item Cross-validation đánh giá độ robust của mô hình.
\end{enumerate}

\begin{ghinho}
\textbf{Lưu ý:} Nội dung này là placeholder tạm thời. Sẽ được bổ sung chi tiết và ví dụ cụ thể sau.
\end{ghinho}
